\section{Background of the study}

The stock market plays an important role in economic development by promoting
capital formation and fostering economic growth. Trading of securities facilitates
savers and users of capital by fund pooling, risk sharing, and transferring wealth
\citep{mosley2008}.

In Nepal, the organized stock market began with the conversion of the Securities
Exchange Centre into Nepal Stock Exchange (NEPSE) Limited in 1993, with floor
trading commencing in 1994. NEPSE remains the sole stock exchange in the country
and has grown substantially in both market capitalization and the number of listed
securities.

Returns from equity investments vary with firm performance and stock price
movements driven by technical, fundamental and market-sentiment factors
\citep{bhattarai2014}. In the Nepalese banking sector, earnings per share (EPS),
price-earnings (P/E) ratio, dividend per share (DPS), book value per share (BVPS)
and return on assets (ROA) are established firm-specific determinants of share
prices \citep{bhattarai2014,pradhan2016,sharma2011}.

The Nepalese banking sector has undergone substantial restructuring through the
merger and acquisition wave of 2013--2020 and the consolidation of early 2023:
Global IME Bank absorbed Bank of Kathmandu, Kumari Bank absorbed Nepal Credit
and Commerce Bank, Nepal Investment Bank merged with Mega Bank to form Nepal
Investment Mega Bank, and Himalayan Bank acquired Civil Bank
\citep{adhikari2023}. Across the broader financial sector, 239 banking and
financial institutions entered merger processes and 177 ceased to exist in their
own names, with the commercial-bank count falling from 32 in 2014 to 20 by
mid-2023. Sector-wide market price per share fell by roughly 46\% in the
periods following the mergers while earnings and dividend per share were
statistically unchanged \citep{adhikari2023,adhikari2023cbr}. The existing
share-price literature is based on pre-consolidation data and samples that no
longer reflect the current structure of the industry. The conflicting findings
on ROA and BVPS across recent studies further underscore the need for a
comprehensive empirical examination of the fundamental determinants of share
prices in the post-consolidation banking sector.
